\Needspace{13\baselineskip}
\Article{18}

Die Methode des vorigen Artikels lieferte die absoluten Werte der Funktionen $T$, $U$ und gab die Bedingungen an, unter denen ihnen gleiche oder entgegengesetzte Vorzeichen zu geben sind; die Vorzeichen selbst sind damit aber noch nicht bestimmt.\quad Für den Fall $k=1$ werden wir dies folgendermaßen ergänzen.

Setzen wir $\rho=\cos\frac12\omega+i\sin\frac12\omega$, sodass $r=\rho\rho$ ist. Wegen $\rho^n=-1$ lässt sich die Gleichung $[8]$ offenbar folgendermaßen schreiben:

\[
\begin{aligned}
W&=2(1+\rho^{n-2})(1+\rho^{-4})(1+\rho^{n-6})(1+\rho^{-8})\\&\qquad\ldots(1+\rho^{-n+4})(1+\rho^2)
\end{aligned}
\]

oder, wenn wir die Faktoren anders anordnen,

\[
\begin{aligned}
W&=2(1+\rho^2)(1+\rho^{-4})(1+\rho^6)(1+\rho^{-8})\\&\qquad\ldots(1+\rho^{-n+4})(1+\rho^{n-2})
\end{aligned}
\]

Nun ist

\[
\begin{aligned}
1+\rho^2&=2\rho\cos\frac12\omega\\
1+\rho^{-4}&=2\rho^{-2}\cos\omega\\
1+\rho^{-6}&=2\rho^3\cos\frac32\omega\\
1+\rho^{-8}&=2\rho^{-4}\cos2\omega,\;\text{usw.}
\end{aligned}
\]

bis hin zu

\[
\begin{aligned}
1+\rho^{-n+4}&=2\rho^{-\frac12n+2}\mathord{,}\cos(\frac14n-1)\omega\\
1+\rho^{n-2}&=2\rho^{\frac12n-1}\cos(\frac14n-\frac12)\omega
\end{aligned}
\]

Somit ergibt sich

\SourceLeaf{41}

\[
\begin{aligned}
W&=2^{\frac12n}\rho^{\frac14n}\cos\frac12\omega\cos\omega\,\frac32\omega\\&\qquad\dotsfour\cos(\frac14n-\frac12)\omega
\end{aligned}
\]

\Needspace{8\baselineskip}
Die in dieses Produkt eingehenden Kosinuswerte sind offenbar sämtlich positiv, der Faktor $\rho^{\frac14n}$ dagegen ist $=\cos45^0+i\sin45^0=(1+i)\surd\frac12$.\quad Daraus schließen wir, dass $W$ das Produkt aus $1+i$ und einer positiven reellen Größe ist. Daher gilt notwendig

\[
\begin{aligned}
W=(1+i)\surd n,\quad T=+\surd n,\quad U=+\surd n
\end{aligned}
\]

\Needspace{20\baselineskip}
\Article{19}

Es wird sich lohnen, sämtliche bisher entwickelten Summationen hier in einer Übersicht zusammenzustellen.\quad Allgemein gilt nämlich

\[
\begin{array}{c|c|c}
T=&U=&\text{für }n\text{ der Form}\\\hline
\pm\surd n&\pm\surd n&4\mu\\
\pm\surd n&0&4\mu+1\\
0&0&4\mu+2\\
0&\pm\surd n&4\mu+3
\end{array}
\]

und wenn wir $k$ $=1$ setzen, muss der Wurzelausdruck das positive Vorzeichen erhalten.\quad Damit ist nun mit voller Strenge bewiesen, was wir in Art. 3 für Primzahlwerte von $n$ durch Induktion bemerkt hatten. Es bleibt nur noch zu zeigen, wie sich die Vorzeichen in allen Fällen für beliebige Werte von $k$ bestimmen lassen.\quad Bevor wir diese Aufgabe jedoch in ihrer ganzen Allgemeinheit angehen können, müssen wir zunächst die Fälle näher betrachten, in denen $n$ eine Primzahl oder eine Primzahlpotenz ist.

\Needspace{14\baselineskip}
\Article{20}

Zunächst sei $n$ eine ungerade Primzahl. Nach den Ausführungen in Art. 10 ist offenbar $W=1+2\Sigma r^a=1+2\Sigma R^{ak}$, wenn $R=\cos\omega+i\sin\omega$ gesetzt wird und $a$, ebenso wie dort, unbestimmt sämtliche quadratischen Reste von $n$ zwischen $1$ und $n-1$ bezeichnet.\quad Bezeichnen wir ebenso durch $b$ unbestimmt sämtliche quadratischen Nichtreste zwischen denselben Grenzen, so sieht man ohne Mühe, dass sämtliche Zahlen $ak$ modulo $n$ entweder zu sämtlichen $a$ oder zu sämtlichen $b$ kongruent sind (ohne Rücksicht auf ihre Reihenfolge), je nachdem, ob $k$ ein Rest oder ein Nichtrest ist.\quad Im ersten Fall ist daher

\SourceLeaf{42}

\[
\begin{aligned}
W=1+2\Sigma R^a=1+R+R^4+R^9+\text{usw.}+R^{(n-1)^2}
\end{aligned}
\]

und somit $W=+\surd n$, wenn $n$ die Form $4\mu+1$ hat, sowie $W=+i\surd n$, wenn $n$ die Form $4\mu+3$ hat.

Im anderen Fall dagegen, in dem $k$ ein Nichtrest von $n$ ist, gilt

\[
\begin{aligned}
W=1+2\Sigma R^b
\end{aligned}
\]

\Needspace{12\baselineskip}
Da nun sämtliche $a$, $b$ offenbar die gesamte Menge der Zahlen $1$, $2$, $3$ $\ldots$ ausfüllen und folglich

\[
\begin{aligned}
\Sigma R^a+\Sigma R^b=R+R^2+R^3+\text{usw.}+R^{n-1}=-1
\end{aligned}
\]

ist, ergibt sich

\[
\begin{aligned}
W=-1-2\Sigma R^a=-(1+R+R^4+R^9+\text{usw.}+R^{(n-1)^2})
\end{aligned}
\]

und somit $W=-\surd n$, wenn $n$ die Form $4\mu+1$ hat, sowie $W=-i\surd n$, wenn $n$ die Form $4\mu+3$ hat.

\Needspace{20\baselineskip}
Daraus folgt also:

\emph{Erstens}, wenn $n$ die Form $4\mu+1$ hat und $k$ ein quadratischer Rest von $n$ ist,

\[
\begin{aligned}
T=+\surd n,\quad U=0
\end{aligned}
\]

\emph{zweitens}, wenn $n$ die Form $4\mu+1$ hat und $k$ ein Nichtrest von $n$ ist,

\[
\begin{aligned}
T=-\surd n,\quad U=0
\end{aligned}
\]

\emph{drittens}, wenn $n$ die Form $4\mu+3$ hat und $k$ ein Rest von $n$ ist,

\[
\begin{aligned}
T=0,\quad U=+\surd n
\end{aligned}
\]

\emph{viertens}, wenn $n$ die Form $4\mu+3$ hat und $k$ ein Nichtrest von $n$ ist,

\[
\begin{aligned}
T=0,\quad U=-\surd n
\end{aligned}
\]
